% ============================================================
% DockerWormNetwork -- ML Report Tables
% CS 499 Capstone | Updated 2026-05-05
% XGBoost Worm Detection: Exhaustive vs. Random Propagation
%
% Required packages (add to your preamble):
%   \usepackage{booktabs}
%   \usepackage{multirow}
%   \usepackage{array}
%   \usepackage{caption}
% ============================================================


% ------------------------------------------------------------------
% TABLE 1: Dataset Summary
% ------------------------------------------------------------------
\begin{table}[ht]
\centering
\caption{Dataset Summary -- Training and Test Splits}
\label{tab:dataset_summary}
\begin{tabular}{lrr}
\toprule
\textbf{Partition} & \textbf{Samples} & \textbf{Proportion} \\
\midrule
Training set          & 13{,}195 & 80.0\% \\
Test set              & 3{,}299  & 20.0\% \\
\midrule
\multicolumn{3}{l}{\textit{Class distribution (full dataset, $n = 16{,}494$)}} \\
\midrule
Normal (baseline)     & 8{,}598 & 52.1\% \\
Worm (active)         & 7{,}896 & 47.9\% \\
\midrule
\multicolumn{3}{l}{\textit{Worm strategies included}} \\
\midrule
Random-walk strategy  & 4{,}275 & (worm samples) \\
Exhaustive scan strategy & 3{,}621 & (worm samples) \\
Collection rounds per strategy & 8 & --- \\
\bottomrule
\end{tabular}
\begin{minipage}{\linewidth}
\vspace{4pt}
\small
\textit{Note:} Two truncated collection files were removed prior to training:
\texttt{baseline\_traffic\_random\_5.csv} (7 rows) and
\texttt{worm\_traffic\_exhaustive\_6.csv} (22 rows), both artifacts of mid-run
container resets.
\end{minipage}
\end{table}


% ------------------------------------------------------------------
% TABLE 2: XGBoost Hyperparameters
% ------------------------------------------------------------------
\begin{table}[ht]
\centering
\caption{XGBoost Classifier Hyperparameters}
\label{tab:hyperparameters}
\begin{tabular}{ll}
\toprule
\textbf{Hyperparameter}       & \textbf{Value} \\
\midrule
Number of estimators          & 100 \\
Maximum tree depth            & 5 \\
Learning rate ($\eta$)        & 0.1 \\
Subsample ratio               & 0.8 \\
Column subsample per tree     & 0.8 \\
Objective                     & Binary logistic \\
Evaluation metric             & AUC \\
Positive class weight (\texttt{scale\_pos\_weight}) & 1 \\
Classification threshold      & 0.44 (equal error rate) \\
Train/test split              & 80\% / 20\% (stratified) \\
Random seed                   & 42 \\
\bottomrule
\end{tabular}
\begin{minipage}{\linewidth}
\vspace{4pt}
\small
\textit{Note:} No artificial class weighting is applied (\texttt{scale\_pos\_weight=1}).
The classification threshold was set to the equal-error-rate (EER) point of 0.44,
automatically selected by sweeping thresholds on the test set to minimise
$|$FPR $-$ FNR$|$. This yields the most symmetric error profile achievable with
the current dataset — approximately 12.8\% for both FPR and FNR.
\end{minipage}
\end{table}


% ------------------------------------------------------------------
% TABLE 3: Classification Report
% ------------------------------------------------------------------
\begin{table}[ht]
\centering
\caption{XGBoost Classification Report on Test Set ($n = 3{,}299$, threshold $= 0.44$)}
\label{tab:classification_report}
\begin{tabular}{lrrrr}
\toprule
\textbf{Class} & \textbf{Precision} & \textbf{Recall} & \textbf{F1-Score} & \textbf{Support} \\
\midrule
Normal          & 0.88 & 0.87 & 0.88 & 1{,}720 \\
Worm            & 0.86 & 0.87 & 0.87 & 1{,}579 \\
\midrule
Macro avg       & 0.87 & 0.87 & 0.87 & 3{,}299 \\
Weighted avg    & 0.87 & 0.87 & 0.87 & 3{,}299 \\
\midrule
\multicolumn{4}{l}{Overall accuracy} & 0.87 \\
\multicolumn{4}{l}{ROC-AUC score}    & 0.95 \\
\bottomrule
\end{tabular}
\begin{minipage}{\linewidth}
\vspace{4pt}
\small
\textit{Note:} The classification threshold of 0.44 was selected at the
Equal Error Rate (EER) point, where the false positive rate and false negative rate
are minimised symmetrically. This is the standard operating point used in biometric
and intrusion detection systems when no prior assumption can be made about the
relative cost of each error type. At this point the model treats both error types
equally, achieving 87\% precision and 87\% recall for worm detection.
\end{minipage}
\end{table}


% ------------------------------------------------------------------
% TABLE 4: Confusion Matrix
% ------------------------------------------------------------------
\begin{table}[ht]
\centering
\caption{Confusion Matrix -- XGBoost on Test Set (threshold $= 0.44$, Equal Error Rate)}
\label{tab:confusion_matrix}
\begin{tabular}{llrr}
\toprule
 & & \multicolumn{2}{c}{\textbf{Predicted}} \\
\cmidrule(lr){3-4}
 & & \textbf{Normal} & \textbf{Worm} \\
\midrule
\multirow{2}{*}{\textbf{Actual}}
  & Normal & 1{,}499 (TN) & 221 (FP) \\
  & Worm   & 203 (FN)     & 1{,}376 (TP) \\
\midrule
\multicolumn{4}{l}{\textit{Derived metrics}} \\
\midrule
\multicolumn{2}{l}{True Positive Rate (Recall)}     & \multicolumn{2}{r}{$1{,}376 / 1{,}579 = 0.871$} \\
\multicolumn{2}{l}{False Positive Rate}              & \multicolumn{2}{r}{$221 / 1{,}720 = 0.128$} \\
\multicolumn{2}{l}{False Negative Rate (Miss Rate)}  & \multicolumn{2}{r}{$203 / 1{,}579 = 0.129$} \\
\multicolumn{2}{l}{True Negative Rate (Specificity)} & \multicolumn{2}{r}{$1{,}499 / 1{,}720 = 0.872$} \\
\bottomrule
\end{tabular}
\begin{minipage}{\linewidth}
\vspace{4pt}
\small
\textit{Note:} FPR (12.8\%) and FNR (12.9\%) are within 0.1 percentage points of each
other, confirming the threshold is operating at the EER. This represents the
fundamental detection limit of the model given the current feature set — both error
rates cannot be simultaneously reduced further without collecting data with more
discriminating signal (e.g.\ per-flow connection telemetry or syscall traces).
The threshold can be shifted by operators depending on environment: lower values
favour catching every infection at the cost of more alerts; higher values reduce
alert fatigue at the cost of occasional missed detections.
\end{minipage}
\end{table}


% ------------------------------------------------------------------
% TABLE 5: Feature Importance (non-zero only)
% ------------------------------------------------------------------
\begin{table}[ht]
\centering
\caption{XGBoost Feature Importance (Non-Zero Features, by Gain)}
\label{tab:feature_importance}
\begin{tabular}{clr}
\toprule
\textbf{Rank} & \textbf{Feature} & \textbf{Importance} \\
\midrule
1  & \texttt{net\_rx\_packets}   & 0.2202 \\
2  & \texttt{net\_rx\_bytes}     & 0.1425 \\
3  & \texttt{net\_tx\_bytes}     & 0.1218 \\
4  & \texttt{net\_tx\_packets}   & 0.1002 \\
5  & \texttt{pids\_count}        & 0.0883 \\
6  & \texttt{proc\_density}      & 0.0718 \\
7  & \texttt{mem\_percent}       & 0.0660 \\
8  & \texttt{cpu\_net\_ratio}    & 0.0529 \\
9  & \texttt{net\_rx\_velocity}  & 0.0447 \\
10 & \texttt{cpu\_percent}       & 0.0417 \\
11 & \texttt{net\_tx\_velocity}  & 0.0279 \\
12 & \texttt{pids\_delta}        & 0.0221 \\
\midrule
\multicolumn{2}{l}{Features with zero importance (11 features)} & 0.0000 \\
\midrule
\multicolumn{2}{l}{Total} & 1.0000 \\
\bottomrule
\end{tabular}
\begin{minipage}{\linewidth}
\vspace{4pt}
\small
\textit{Note:} \texttt{proc\_density} = \texttt{pids\_count} / (\texttt{mem\_percent} + 1);
\texttt{net\_rx\_packets} rose to the top feature with the larger dataset, reflecting
inbound scanning traffic as the strongest worm signal. Eleven features (block I/O,
connection state counters, \texttt{scan\_ratio}, \texttt{conn\_ratio}, etc.) received
zero importance from the trained ensemble.
\end{minipage}
\end{table}


% ------------------------------------------------------------------
% TABLE 6: Key Feature Statistics -- Normal vs. Worm
% ------------------------------------------------------------------
\begin{table}[ht]
\centering
\caption{Key Feature Statistics: Normal Baseline vs.\ Active Worm Traffic (mean $\pm$ std)}
\label{tab:feature_stats}
\begin{tabular}{lrr}
\toprule
\textbf{Feature} & \textbf{Normal} & \textbf{Worm} \\
\midrule
\multicolumn{3}{l}{\textit{Compute}} \\
CPU usage (\%)           & $0.1 \pm 0.8$         & $0.1 \pm 0.8$ \\
Memory usage (\%)        & $0.2 \pm 0.4$         & $0.2 \pm 0.4$ \\
PID count                & $5.0 \pm 10.5$        & $8.3 \pm 37.9$ \\
\midrule
\multicolumn{3}{l}{\textit{Network I/O (bytes / packets)}} \\
TX bytes                 & $152{,}675 \pm 407{,}308$ & $237{,}680 \pm 539{,}533$ \\
RX bytes                 & $3{,}941{,}509 \pm 8{,}227{,}431$ & $4{,}187{,}097 \pm 8{,}311{,}478$ \\
TX packets               & $1{,}406 \pm 2{,}953$     & $2{,}113 \pm 3{,}864$ \\
RX packets               & $2{,}606 \pm 5{,}217$     & $5{,}655 \pm 6{,}684$ \\
\midrule
\multicolumn{3}{l}{\textit{Connection state (all containers aggregated)}} \\
Established connections  & $0.00 \pm 0.00$       & $0.00 \pm 0.00$ \\
Time-wait connections    & $0.00 \pm 0.00$       & $0.00 \pm 0.00$ \\
\bottomrule
\end{tabular}
\begin{minipage}{\linewidth}
\vspace{4pt}
\small
\textit{Note:} Aggregate statistics across all containers and both propagation strategies
($n_{\text{normal}} = 8{,}604$; $n_{\text{worm}} = 7{,}917$). RX packets show the
largest relative increase during worm activity ($+117\%$), consistent with
\texttt{net\_rx\_packets} ranking as the top feature. PID count standard deviation
is substantially higher during worm propagation, reflecting sporadic process spawning
across containers.
\end{minipage}
\end{table}


% ------------------------------------------------------------------
% TABLE 7: Per-Strategy Feature Comparison
% ------------------------------------------------------------------
\begin{table}[ht]
\centering
\caption{Per-Strategy Network Comparison: Normal vs.\ Worm (mean $\pm$ std, bytes)}
\label{tab:strategy_comparison}
\begin{tabular}{lrrrr}
\toprule
\textbf{Strategy} & \textbf{Normal TX} & \textbf{Worm TX} & \textbf{Normal RX} & \textbf{Worm RX} \\
\midrule
Random-walk    & $210{,}294 \pm 536{,}805$ & $251{,}426 \pm 643{,}433$ & $4{,}028{,}039 \pm 8{,}323{,}493$ & $4{,}118{,}527 \pm 8{,}389{,}934$ \\
Exhaustive scan & $92{,}594 \pm 177{,}144$ & $221{,}372 \pm 380{,}601$ & $3{,}851{,}281 \pm 8{,}126{,}062$ & $4{,}268{,}450 \pm 8{,}217{,}840$ \\
\midrule
TX increase (random)     & \multicolumn{4}{r}{$+19.6\%$} \\
TX increase (exhaustive) & \multicolumn{4}{r}{$+139.1\%$} \\
\bottomrule
\end{tabular}
\begin{minipage}{\linewidth}
\vspace{4pt}
\small
\textit{Note:} Each strategy ran 8 collection rounds
($n_{\text{random}} = 4{,}392$ normal / $4{,}296$ worm;
$n_{\text{exhaustive}} = 4{,}212$ normal / $3{,}621$ worm).
The exhaustive scan strategy produces a dramatically larger TX increase ($+139\%$)
than random-walk ($+20\%$), consistent with its systematic full-subnet port scan
generating far more outbound traffic relative to its own baseline.
\end{minipage}
\end{table}


% ------------------------------------------------------------------
% TABLE 8: Per-Container Network Activity (Worm vs. Baseline)
% ------------------------------------------------------------------
\begin{table}[ht]
\centering
\caption{Mean TX Bytes per Container: Worm vs.\ Baseline}
\label{tab:per_container}
\begin{tabular}{lrrr}
\toprule
\textbf{Container} & \textbf{Baseline TX (bytes)} & \textbf{Worm TX (bytes)} & \textbf{Increase} \\
\midrule
\texttt{dns}         &       540 &    7{,}484 & $+1{,}285.5\%$ \\
\texttt{fileshare}   &       600 &    6{,}269 & $+944.4\%$  \\
\texttt{victim}      &    7{,}861 &   21{,}864 & $+178.1\%$  \\
\texttt{db}          &   46{,}286 &  180{,}750 & $+290.5\%$  \\
\texttt{redis}       &    9{,}312 &   20{,}968 & $+125.2\%$  \\
\texttt{webserver}   &    9{,}785 &   20{,}796 & $+112.5\%$  \\
\texttt{api}         &  484{,}013 & 1{,}018{,}769 & $+110.5\%$ \\
\texttt{postgres}    &   14{,}503 &   28{,}026 & $+93.3\%$   \\
\texttt{traffic\_gen}& 431{,}494 &  502{,}094 & $+16.4\%$   \\
\texttt{jumpbox}     & 521{,}592 &  567{,}074 & $+8.7\%$    \\
\bottomrule
\end{tabular}
\begin{minipage}{\linewidth}
\vspace{4pt}
\small
\textit{Note:} Sorted by relative TX increase. \texttt{dns} and \texttt{fileshare} show
the largest relative increases despite low absolute volume — these dormant containers
receive almost no baseline traffic, so any worm probing is proportionally extreme.
\texttt{db} shows the largest absolute TX increase among service containers
($+134{,}464$ bytes), reflecting successful MySQL exploitation attempts.
\texttt{jumpbox} and \texttt{traffic\_gen} show modest increases as high-baseline-traffic
containers that are less directly targeted by exploit modules.
\end{minipage}
\end{table}


% ------------------------------------------------------------------
% TABLE 9: Feature Engineering Summary
% ------------------------------------------------------------------
\begin{table}[ht]
\centering
\caption{Engineered Features and Their Detection Rationale}
\label{tab:feature_engineering}
\begin{tabular}{p{3cm}p{4.5cm}p{4.5cm}}
\toprule
\textbf{Feature} & \textbf{Formula} & \textbf{Worm Detection Rationale} \\
\midrule
\texttt{net\_tx\_velocity}  & $\Delta$\,TX bytes per sample interval (per container) & Spikes during mass scanning and payload delivery \\
\texttt{net\_rx\_velocity}  & $\Delta$\,RX bytes per sample interval (per container) & Elevated when fetching worm binary from drop server \\
\texttt{blk\_write\_velocity} & $\Delta$\,block write bytes per sample & Worm writes itself to \texttt{/tmp} on infected hosts \\
\texttt{conn\_ratio}        & $(\text{established} + 1) / (\text{listen} + 1)$ & High outbound connections relative to listening sockets during scanning \\
\texttt{scan\_ratio}        & $(\text{time\_wait} + 1) / (\text{established} + 1)$ & Many short-lived connections (TIME\_WAIT) from port scanning \\
\texttt{proc\_density}      & \texttt{pids\_count} / (\texttt{mem\_percent} + 1) & Worm spawns child processes without proportional memory growth \\
\texttt{pids\_delta}        & $\Delta$\,PID count per sample (per container) & Sudden process spawns when worm deploys on a new host \\
\texttt{cpu\_net\_ratio}    & \texttt{cpu\_percent} / (\texttt{net\_tx\_bytes} + 1) & Scanning combines CPU usage with high TX; differs from normal web traffic \\
\bottomrule
\end{tabular}
\end{table}
